\documentclass[aspectratio=169]{beamer}
\usepackage[T1]{fontenc}
\usepackage[ngerman]{babel}
\usepackage{lmodern}

% Design-Anpassungen 
\usetheme{madrid}
\usecolortheme{default}

% Metadaten
\title[Cybersicherheitskurs]{Entwicklung eines interaktiven Cybersicherheitskurses}
\subtitle{Maturarbeit}
\author{Laurin Ott}
\institute{Kantonsschule Im Lee, Winterthur}
\date{\today}

\begin{document}

% Titelfolie
\begin{frame}
    \titlepage
\end{frame}

% Inhaltsverzeichnis
\begin{frame}{Inhaltsverzeichnis}
    \tableofcontents
\end{frame}

\section{Einleitung \& Motivation}
\begin{frame}{Einleitung \& Motivation}
    \begin{itemize}
        \item \textbf{Das Problem:} Digitale Bedrohungen (Phishing, Social Engineering, unsichere Passwörter) nehmen zu.
        \item \textbf{Die Zielsetzung:} Untersuchung, wie effektiv ein gezielter, interaktiver Kurs das Sicherheitsbewusstsein von Schülerinnen und Schülern verbessern kann.
        \item \textbf{Der Ansatz:} Nicht nur Theorie vermitteln, sondern aktiv und interaktiv erleben lassen.
    \end{itemize}
\end{frame}

\section{Das Konzept}
\begin{frame}{Das Konzept: Die Web-App}
    \begin{itemize}
        \item Entwicklung einer \textbf{eigenen Web-App} als interaktiver Cybersicherheitskurs.
        \item Spielerische und verständliche Vermittlung von typischen Online-Gefahren:
        \begin{itemize}
            \item Phishing-Mails erkennen
            \item Social Engineering (Manipulationstechniken)
            \item Erstellung und Verwaltung sicherer Passwörter
        \end{itemize}
    \end{itemize}
\end{frame}

\section{Methodik \& Evaluation}
\begin{frame}{Methodik: Vor- und Nachtest}
    \begin{itemize}
        \item \textbf{Messung des Lern- und Verhaltenseffekts:}
        \begin{enumerate}
            \item \textbf{Vortest:} Erfassung des Wissensstands vor dem Kurs.
            \item \textbf{Durchführung:} Absolvieren des interaktiven Web-App-Kurses.
            \item \textbf{Nachtest:} Überprüfung, ob sich das Verhalten und Bewusstsein verbessert haben.
        \end{enumerate}
        \item Ermöglicht eine datenbasierte Auswertung der Kurseffektivität.
    \end{itemize}
\end{frame}

\section{Umsetzung \& Integration}
\begin{frame}{Umsetzung an der Kanti Im Lee}
    \begin{itemize}
        \item \textbf{Testphase \& Untersuchung:} 
        \begin{itemize}
            \item Durchführung direkt an der \textbf{Kantonsschule Im Lee in Winterthur}.
        \end{itemize}
        \item \textbf{Technische Umsetzung \& Open Source:}
        \begin{itemize}
            \item Kompletter Code und Projektverlauf offen auf \textbf{GitHub} dokumentiert.
        \end{itemize}
        \item \textbf{Nachhaltige Integration:}
        \begin{itemize}
            \item Einbettung des fertigen Kurses direkt in die \textbf{Moodle-Lernplattform} der Schule für den zukünftigen Unterricht.
        \end{itemize}
    \end{itemize}
\end{frame}

\section{Fazit \& Ausblick}
\begin{frame}{Fazit \& Ausblick}
    \begin{itemize}
        \item \textbf{Mehrwert:} Verbindung von praktischer Softwareentwicklung und pädagogischem Nutzen für den Schulalltag.
        \item \textbf{Ausblick:} Nachhaltiger Einsatz des Kurses an der Kanti Im Lee und potenziell darüber hinaus.
    \end{itemize}
\end{frame}

\end{document}